% ============================================================
% 3.2.3.2 Employment and Unemployment — Worksheet
% AQA AS Economics
% Sources: output/authoritative/as_syllabus.json (3.2.3.2)
%          output/authoritative/sow.json (as > 3.2 > 3.2.3)
% Total marks: 60  |  Suggested time: 2 hours (inc. group discussion)
% ============================================================

\documentclass[12pt,a4paper]{article}
\usepackage[margin=1in,headheight=15pt]{geometry}
\usepackage{enumitem}
\usepackage{amsmath}
\usepackage{fancyhdr}
\usepackage{booktabs}
\usepackage{array}
\usepackage{tikz}
\usepackage{pgfplots}
\pgfplotsset{compat=1.18}
\RequirePackage{tcolorbox}
\tcbuselibrary{skins}

% Key Point — yellow; for the single most important idea on a slide
\newtcolorbox{keypoint}{
    enhanced,
    colback=yellow!12,
    colframe=yellow!60!black,
    fonttitle=\bfseries,
    title=Key Point,
    arc=2pt,
    boxrule=0.8pt,
}

% Formula — blue; for named equations and their variable key
\newtcolorbox{formula}[1][Formula]{
    enhanced,
    colback=blue!5,
    colframe=blue!45!black,
    fonttitle=\bfseries,
    title=#1,
    arc=2pt,
    boxrule=0.8pt,
}

% Example Question — teal; for exam-style questions posed to students
\newtcolorbox{examplequestion}[1][Example Question]{
    enhanced,
    colback=teal!6,
    colframe=teal!55!black,
    fonttitle=\bfseries,
    title=#1,
    arc=2pt,
    boxrule=0.8pt,
}

% Example Solution — green; always follows an examplequestion block
\newtcolorbox{examplesolution}[1][Solution]{
    enhanced,
    colback=green!6,
    colframe=green!45!black,
    fonttitle=\bfseries,
    title=#1,
    arc=2pt,
    boxrule=0.8pt,
}

% ---- Worksheet environments --------------------------------

% Scenario — blue; for group discussion prompts on worksheets
% Usage: \begin{scenario}{Scenario 1: Title with commas, fine} ... \end{scenario}
\newtcolorbox{scenario}[1]{
    enhanced,
    colback=blue!5,
    colframe=blue!45!black,
    fonttitle=\bfseries,
    title={#1},
    arc=2pt,
    boxrule=0.8pt,
}

% Instructions — blue; for worksheet-level instruction blocks
\newtcolorbox{instructions}{
    enhanced,
    colback=blue!5,
    colframe=blue!45!black,
    fonttitle=\bfseries,
    title=Instructions,
    arc=2pt,
    boxrule=0.8pt,
}

% Extract — grey; for data response source material
% Usage: \begin{extract}{Source: Author, Year} ... \end{extract}
\newtcolorbox{extract}[1]{
    enhanced,
    colback=gray!8,
    colframe=gray!50,
    fonttitle=\bfseries,
    title={#1},
    arc=2pt,
    boxrule=0.8pt,
}


\pagestyle{fancy}
\fancyhf{}
\lhead{Dr Jac Thomas}
\rhead{Topic index: 3.2.3.2}
\rfoot{Page \thepage}

\begin{document}

% ---- Title block -------------------------------------------
\begin{center}
    \Large\textbf{Employment and Unemployment}\\[5pt]
    \normalsize Topic index: 3.2.3.2 \quad|\quad AQA AS Economics\\[3pt]
    \textit{Total Marks: 60 \quad Suggested Time: 2 hours (including group discussion)}
\end{center}

\vspace{0.8em}

\begin{instructions}
    Answer all questions in the spaces provided. Show your working clearly for calculation questions. For the group discussion, record your group's agreed conclusions in the spaces given.
\end{instructions}

\vspace{1em}

% ============================================================
% SECTION A: KEY TERMS
% ============================================================
\section*{Section A: Key Terms (10 marks)}

\textit{Define each of the following terms. Each definition is worth 2 marks.}

\vspace{0.5em}

\begin{enumerate}[leftmargin=*, label=\textbf{\arabic*.}]

    \item \textbf{Unemployment} \hfill (2 marks)\\[3cm]

    \item \textbf{The Claimant Count} \hfill (2 marks)\\[3cm]

    \item \textbf{The Labour Force Survey (LFS)} \hfill (2 marks)\\[3cm]

    \item \textbf{Structural unemployment} \hfill (2 marks)\\[3cm]

    \item \textbf{Cyclical unemployment} \hfill (2 marks)\\[3cm]

\end{enumerate}

\newpage

% ============================================================
% SECTION B: GROUP DISCUSSION
% ============================================================
\section*{Section B: Group Discussion (approximately 20 minutes)}

\textit{In your groups, read each scenario. For each, decide:}
\begin{enumerate}[label=(\alph*), topsep=2pt]
    \item What type of unemployment does this represent?
    \item What is the main cause?
    \item What policy would best address it, and why?
\end{enumerate}

\vspace{0.8em}

\begin{scenario}{Scenario 1: Tata Steel, Port Talbot (2024)}
    In September 2024, Tata Steel announced the closure of the last blast furnaces at its Port Talbot steelworks, with approximately 2,800 jobs lost. Many of the workers had spent their entire careers in the steel industry and held highly specialised skills not easily transferable to other sectors.
\end{scenario}

\vspace{0.5em}
\textbf{Your analysis:}\\
Type of unemployment: \dotfill \\[0.4em]
Main cause: \dotfill \\[0.4em]
Appropriate policy: \dotfill \\[0.4em]
Justification: \dotfill

\vspace{1em}

\begin{scenario}{Scenario 2: The 2008 Financial Crisis}
    Following the collapse of Lehman Brothers in 2008, UK banks sharply reduced lending to businesses and consumers. Consumer spending fell, business investment collapsed, and the UK entered recession. Unemployment rose from around 5\% in 2007 to a peak of 8.5\% in 2011.
\end{scenario}

\vspace{0.5em}
\textbf{Your analysis:}\\
Type of unemployment: \dotfill \\[0.4em]
Main cause: \dotfill \\[0.4em]
Appropriate policy: \dotfill \\[0.4em]
Justification: \dotfill

\vspace{1em}

\begin{scenario}{Scenario 3: UK Coastal Tourism}
    Every autumn, seaside hotels, ice cream sellers, and seasonal theme parks lay off a significant proportion of their workforce. Employment in these sectors typically falls sharply between October and March, before picking up again from April onwards.
\end{scenario}

\vspace{0.5em}
\textbf{Your analysis:}\\
Type of unemployment: \dotfill \\[0.4em]
Main cause: \dotfill \\[0.4em]
Is this necessarily a problem for the economy? \dotfill \\[0.4em]
Could or should policy reduce it? \dotfill

\newpage

% ============================================================
% SECTION C: MULTIPLE CHOICE
% ============================================================
\section*{Section C: Multiple Choice (10 marks)}

\textit{Circle the correct answer. Each question is worth 1 mark.}

\vspace{0.5em}

\begin{enumerate}[leftmargin=*, label=\textbf{\arabic*.}]

    \item The claimant count measures unemployment by counting people who are:
    \begin{enumerate}[label=(\Alph*), topsep=2pt]
        \item Without work and actively seeking employment
        \item Economically inactive
        \item Claiming Universal Credit or Jobseeker's Allowance
        \item Out of work for more than 6 months
    \end{enumerate}

    \vspace{0.4em}

    \item Under the Labour Force Survey (ILO) definition, a person is classified as unemployed if they are:
    \begin{enumerate}[label=(\Alph*), topsep=2pt]
        \item Not currently in paid employment
        \item Without work, available for work, and actively seeking work
        \item Claiming state benefits
        \item Working fewer than 16 hours per week
    \end{enumerate}

    \vspace{0.4em}

    \item A car manufacturer replaces 600 assembly workers with robots. This best illustrates:
    \begin{enumerate}[label=(\Alph*), topsep=2pt]
        \item Frictional unemployment
        \item Seasonal unemployment
        \item Structural unemployment
        \item Cyclical unemployment
    \end{enumerate}

    \vspace{0.4em}

    \item A recent graduate spends three months searching for their first job after university. This most likely represents:
    \begin{enumerate}[label=(\Alph*), topsep=2pt]
        \item Structural unemployment
        \item Seasonal unemployment
        \item Frictional unemployment
        \item Cyclical unemployment
    \end{enumerate}

    \vspace{0.4em}

    \item Which of the following is a demand-side cause of rising unemployment?
    \begin{enumerate}[label=(\Alph*), topsep=2pt]
        \item Workers lacking relevant skills for available jobs
        \item A rise in income tax reducing consumer spending and aggregate demand
        \item Restrictive employment protection legislation reducing labour market flexibility
        \item A geographical mismatch between workers and job vacancies
    \end{enumerate}

    \vspace{0.4em}

    \item A depreciation of the pound sterling is most likely to reduce UK unemployment by:
    \begin{enumerate}[label=(\Alph*), topsep=2pt]
        \item Lowering import costs for UK firms
        \item Making UK exports cheaper and more competitive abroad
        \item Increasing structural unemployment in export industries
        \item Reducing inflationary pressure
    \end{enumerate}

    \vspace{0.4em}

    \item An economy experiencing a negative output gap is most likely to show:
    \begin{enumerate}[label=(\Alph*), topsep=2pt]
        \item Inflationary pressure and falling unemployment
        \item Spare productive capacity and above-average unemployment
        \item Output at its full-employment potential
        \item An increase in long-run aggregate supply
    \end{enumerate}

    \vspace{0.4em}

    \item Which policy is most appropriate for reducing structural unemployment?
    \begin{enumerate}[label=(\Alph*), topsep=2pt]
        \item Cutting interest rates to boost consumer spending
        \item Increasing government spending on infrastructure projects
        \item Introducing retraining schemes to improve occupational mobility
        \item Reducing income tax to stimulate aggregate demand
    \end{enumerate}

    \vspace{0.4em}

    \item An economic slowdown in the Eurozone would most directly affect UK unemployment through:
    \begin{enumerate}[label=(\Alph*), topsep=2pt]
        \item Lower UK export demand, reducing output and employment in export industries
        \item An appreciation of the pound sterling raising export prices
        \item Higher UK import prices reducing real household incomes
        \item A fall in UK interest rates reducing business investment
    \end{enumerate}

    \vspace{0.4em}

    \item The claimant count tends to give a lower unemployment figure than the LFS because:
    \begin{enumerate}[label=(\Alph*), topsep=2pt]
        \item The LFS only surveys workers in large firms
        \item Not all unemployed people are eligible for or choose to claim benefits
        \item The claimant count includes economically inactive people
        \item The LFS counts only those unemployed for more than 6 months
    \end{enumerate}

\end{enumerate}

\newpage

% ============================================================
% SECTION D: QUANTITATIVE SKILLS
% ============================================================
\section*{Section D: Quantitative Skills (8 marks)}

\textit{Show all working clearly. State the formula before substituting values.}

\vspace{0.8em}

The table below shows UK labour market data for two consecutive years.

\vspace{0.5em}

\begin{table}[h]
    \centering
    \renewcommand{\arraystretch}{1.3}
    \begin{tabular}{@{}lrr@{}}
        \toprule
        \textbf{Measure} & \textbf{2023} & \textbf{2024} \\ \midrule
        Labour force (millions)       & 35.00  & 35.00  \\
        Number employed (millions)    & 33.60  & 33.25  \\
        Number unemployed -- LFS (millions) & 1.40 & ?  \\
        Claimant count (millions)     & 0.98   & ---    \\ \bottomrule
    \end{tabular}
\end{table}

\vspace{0.8em}

\begin{enumerate}[leftmargin=*, label=\textbf{\arabic*.}]

    \item Calculate the LFS unemployment rate for 2023. \hfill (2 marks)\\[0.2em]
    \textit{Show your working:}
    \vspace{3cm}

    \item Calculate the number of people unemployed (LFS measure) in 2024. \hfill (1 mark)\\[0.2em]
    \textit{Show your working:}
    \vspace{2cm}

    \item Using your answer to question 2, calculate the LFS unemployment rate for 2024. \hfill (2 marks)\\[0.2em]
    \textit{Show your working:}
    \vspace{3cm}

    \item Calculate the percentage change in the number of people unemployed between 2023 and 2024. \hfill (2 marks)\\[0.2em]
    \textit{Show your working:}
    \vspace{3cm}

    \item In 2023, express the claimant count as a percentage of the LFS unemployment figure. \hfill (1 mark)\\[0.2em]
    \textit{Show your working:}
    \vspace{2cm}

\end{enumerate}

\newpage

% ============================================================
% SECTION E: DIAGRAMS
% ============================================================
\section*{Section E: Diagrams (10 marks)}

\begin{enumerate}[leftmargin=*, label=\textbf{\arabic*.}]

    \item Draw a fully labelled AD/AS diagram to illustrate an economy experiencing cyclical unemployment. Your diagram should clearly show the negative output gap between actual output and full-employment output.\hfill (5 marks)

    \vspace{0.4em}
    \begin{center}
        \begin{tikzpicture}
            \begin{axis}[
                axis lines = middle,
                xlabel = {Real GDP},
                ylabel = {Price Level},
                xmin=0, xmax=10, ymin=0, ymax=10,
                xtick=\empty, ytick=\empty,
                width=10.5cm, height=7cm,
                xlabel style={at={(axis description cs:1,0)}, anchor=north east, yshift=-10pt, font=\small},
                ylabel style={at={(axis description cs:0,1)}, anchor=south east, font=\small},
            ]
            % Student draws here
            \end{axis}
        \end{tikzpicture}
    \end{center}

    \vspace{0.6em}

    \item On the diagram below, show how a rise in aggregate demand could reduce cyclical unemployment. Start from the position in question 1. Label the initial equilibrium ($E_1$) and the new equilibrium ($E_2$).\hfill (5 marks)

    \vspace{0.4em}
    \begin{center}
        \begin{tikzpicture}
            \begin{axis}[
                axis lines = middle,
                xlabel = {Real GDP},
                ylabel = {Price Level},
                xmin=0, xmax=10, ymin=0, ymax=10,
                xtick=\empty, ytick=\empty,
                width=10.5cm, height=7cm,
                xlabel style={at={(axis description cs:1,0)}, anchor=north east, yshift=-10pt, font=\small},
                ylabel style={at={(axis description cs:0,1)}, anchor=south east, font=\small},
            ]
            % Student draws here
            \end{axis}
        \end{tikzpicture}
    \end{center}

\end{enumerate}

\newpage

% ============================================================
% SECTION F: SHORT RESPONSE
% ============================================================
\section*{Section F: Short Response (7 marks)}

\begin{enumerate}[leftmargin=*, label=\textbf{\arabic*.}]

    \item Explain \textbf{two} differences between the claimant count and the Labour Force Survey as measures of unemployment. \hfill (4 marks)\\[6.5cm]

    \item Explain one supply-side factor that could cause a rise in unemployment in the UK. \hfill (3 marks)\\[5.5cm]

\end{enumerate}

\newpage

% ============================================================
% SECTION G: ESSAY
% ============================================================
\section*{Section G: Essay (15 marks)}

\begin{enumerate}[leftmargin=*, label=\textbf{\arabic*.}]

    \item \textquotedblleft To what extent do the appropriate policies for reducing unemployment depend on the type of unemployment?\textquotedblright \hfill (15 marks)

    \vspace{0.8em}

    \begin{keypoint}
        \textbf{Evaluation tips:} Consider what each of the four types of unemployment has in common and where they differ. Think about whether demand-side policies (e.g.\ expansionary fiscal or monetary policy) can address all types equally well. Consider supply-side policies such as retraining programmes, improving geographical and occupational mobility, and reducing information failures in job markets. Can one policy simultaneously reduce multiple types? Think about time lags before policies take effect, and any trade-offs with other macroeconomic objectives (e.g.\ inflation).
    \end{keypoint}

    \vspace{14cm}

\end{enumerate}

\vfill
\begin{center}
    \textbf{[Total: 60 marks]}
\end{center}

\end{document}
